\documentclass{article}
\usepackage[T1]{fontenc}
\usepackage{lmodern}
\usepackage{graphicx}
\usepackage{amsmath,amssymb}
\usepackage[a4paper,margin=2cm]{geometry}
\usepackage[absolute,overlay]{textpos}
\setlength{\TPHorizModule}{1mm}
\setlength{\TPVertModule}{1mm}
\usepackage[indonesian]{babel}
\usepackage{hyperref}
\setlength{\emergencystretch}{2em}
% unit: Halaman 32

\begin{document}

% === Halaman 32 (python/native) ===

Penggandaan Titik di dalam EC pada GF(p)

Misalkan P(xp, yp) yang dalam hal ini yp  0.  

Penggandaan titik:  2P = R

Koordinat Titik R: 

     xr  m2 – 2xp  (mod p)

     yr  m(xp – xr) – yp (mod p)

Yang dalam hal ini,

Jika yp = 0 maka m tidak terdefinisi sehingga 2P = O

32

𝑚 3𝑥𝑝

2 + 𝑎

2𝑦𝑝

 (mod 𝑝)

\end{document}
